\documentclass[11pt]{article}
\usepackage[letterpaper,margin=1in]{geometry}
\usepackage{amsmath,amssymb,amsthm,mathtools}
\usepackage{graphicx}
\usepackage{microtype}
\usepackage{xcolor}
\usepackage{enumitem}
\usepackage[hidelinks]{hyperref}

\newtheorem{theorem}{Theorem}
\newtheorem{proposition}{Proposition}
\newcommand{\R}{\mathbb R}
\newcommand{\E}{\mathbb E}
\newcommand{\g}{\gamma}
\newcommand{\T}{\mathcal T}
\newcommand{\Pp}{\mathcal P}
\newcommand{\bmo}{\mathrm{bmo}}
\newcommand{\BMO}{\mathrm{BMO}}
\newcommand{\Ga}{\Gamma}

\title{Gaussian drift separates the intrinsic semigroup BMO norms}
\author{Automated mathematics run \texttt{fa\_banach\_001}\\
Lane 10 --- candidate counterexample for expert review}
\date{9 August 2026}

\begin{document}
\maketitle

\begin{center}
\fbox{\parbox{0.91\textwidth}{
\textbf{Status.} Candidate counterexample, likely valid.  For the classical
one-dimensional Ornstein--Uhlenbeck semigroup, bounded coordinate truncations
have uniformly bounded $\bmo(\T)$ norm while their $\BMO(\T)$ norms and
$\sup_t\|T_tf-T_{2t}f\|_\infty$ diverge.  The coordinate also separates
$\BMO(\Ga)$ from $\BMO(\Pp)$ and $\BMO(\partial)$.  This answers the universal
readings of open problem (ii)(a),(c),(c') of Junge--Mei negatively, but does
not characterize all semigroups for which the comparisons hold.
}}
\end{center}

\section{Source question and definitions}

Let $\T=(T_t)_{t\geq0}$ be a standard Markov semigroup.  In the commutative
case, the two norms in open problem (ii)(c) of Junge and Mei
\cite{JungeMei} are
\begin{align}
 \|f\|_{\bmo(\T)}^2
  &=\sup_{t>0}\big\|T_t|f|^2-|T_tf|^2\big\|_\infty,\label{eq:bmo}\\
 \|f\|_{\BMO(\T)}^2
  &=\sup_{t>0}\big\|T_t|f-T_tf|^2\big\|_\infty.\label{eq:BMO}
\end{align}
They ask for conditions under which these norms are equivalent, and single
out the estimate
\begin{equation}\label{eq:cprime}
 \sup_{t>0}\|T_tf-T_{2t}f\|_\infty
 \leq C\|f\|_{\bmo(\T)}.
\end{equation}
Part (a) asks for comparison of the Poisson-semigroup norm $\BMO(\Pp)$,
the vertical square-function norm $\BMO(\partial)$, and the carré-du-champ
norm $\BMO(\Ga)$.

\begin{figure}[ht]
  \centering
  \includegraphics[width=0.98\textwidth]{figures/open_problem_crop.png}
  \caption{Open problem (ii) on printed page 33 of the source paper
  \cite{JungeMei}.}
\end{figure}

We prove that even the strongest standard regularity package does not imply
(a), (c), or (c').  Work on the probability space
\[
 d\g(x)=(2\pi)^{-1/2}e^{-x^2/2}\,dx
\]
with the Ornstein--Uhlenbeck semigroup
\begin{equation}\label{eq:mehler}
 T_tf(x)=\E f\!\left(e^{-t}x+\sqrt{1-e^{-2t}}\,Z\right),
 \qquad Z\sim N(0,1).
\end{equation}
It is a symmetric, ergodic, analytic and hypercontractive Markov semigroup.
For its positive generator $A=-\partial_x^2+x\partial_x$,
$\Ga(f)=|f'|^2$ and
\[
 \Ga_2(f)=|f''|^2+|f'|^2\geq \Ga(f).
\]
Thus the example satisfies strictly more than the $\Ga_2\geq0$ hypothesis
used in the source.

\section{Proof intuition}

Mehler's formula turns the expression in \eqref{eq:bmo} into the conditional
variance of $f$ under a Gaussian transition kernel.  The truncations
$f_R(x)=\max(-R,\min(x,R))$ are $1$-Lipschitz, so their transition variances
never exceed the variance $1-e^{-2t}$ of the Gaussian input.  Their
$\bmo(\T)$ norms are therefore uniformly bounded.

The other norm sees something that conditional variance misses: the OU drift
of the mean.  On the untruncated coordinate $h(x)=x$, the mean changes from
$e^{-t}x$ to $e^{-2t}x$, an unbounded function of $x$.  Since $f_R\to h$ in
$L_2(\g)$, the same mean-drift term for $f_R$ must have unbounded
$L_\infty$ norm as $R\to\infty$.  Kadison's inequality puts this drift term
below the $\BMO(\T)$ norm.  This proves both failures using bounded functions.

For the subordinated Poisson semigroup, the carré du champ sees only the
bounded derivative of the coordinate, while the oscillation and vertical
derivative norms still see its unbounded drift.  That is the mechanism behind
the separation in part (a).

\section{The bounded counterexample to (c) and (c')}

For $R>0$, set
\begin{equation}\label{eq:trunc}
 f_R(x)=(-R)\vee x\wedge R.
\end{equation}

\begin{theorem}\label{thm:main}
For the Ornstein--Uhlenbeck semigroup \eqref{eq:mehler}:
\begin{enumerate}[label=\textup{(\roman*)},leftmargin=2.3em]
\item $\|f_R\|_{\bmo(\T)}\leq1$ for every $R>0$;
\item for every fixed $t>0$,
\[
 \|(T_t-T_{2t})f_R\|_\infty\longrightarrow\infty
 \qquad(R\to\infty);
\]
\item consequently $\|f_R\|_{\BMO(\T)}\to\infty$.
\end{enumerate}
Hence \eqref{eq:cprime} is false and the two norms in open problem (ii)(c)
are not equivalent, even on bounded functions for one fixed classical
semigroup satisfying $\Ga_2\geq\Ga$.
\end{theorem}

\begin{proof}
Fix $t>0$ and $x\in\R$, and let
\[
 Y=e^{-t}x+\sqrt{1-e^{-2t}}\,Z.
\]
The quantity in \eqref{eq:bmo} is
\[
 T_tf_R^2(x)-(T_tf_R(x))^2=\operatorname{Var}(f_R(Y)).
\]
If $Y'$ is an independent copy of $Y$, then the $1$-Lipschitz property gives
\[
 \operatorname{Var}(f_R(Y))
 =\frac12\E|f_R(Y)-f_R(Y')|^2
 \leq\frac12\E|Y-Y'|^2
 =1-e^{-2t}\leq1.
\]
Taking the two suprema proves (i).

Let $h(x)=x$.  Since $f_R\to h$ in $L_2(\g)$ and every $T_s$ is an
$L_2$ contraction,
\begin{equation}\label{eq:L2limit}
 (T_t-T_{2t})f_R
 \longrightarrow (e^{-t}-e^{-2t})h
 \quad\hbox{in }L_2(\g).
\end{equation}
The limit is essentially unbounded.  If the $L_\infty$ norms on the left did
not tend to infinity, some subsequence would be bounded in $L_\infty$ by one
constant.  An $L_2$-convergent sub-subsequence converges almost everywhere,
forcing the limit in \eqref{eq:L2limit} to have the same $L_\infty$ bound, a
contradiction.  This proves (ii).

Finally, positivity and Kadison's inequality give, pointwise,
\begin{align*}
 |T_tf_R-T_{2t}f_R|^2
 &=|T_t(f_R-T_tf_R)|^2\\
 &\leq T_t|f_R-T_tf_R|^2.
\end{align*}
Thus $\|(T_t-T_{2t})f_R\|_\infty\leq\|f_R\|_{\BMO(\T)}$, and (iii)
follows.
\end{proof}

The separation also occurs at the level of the natural $L_2$ spaces, without
truncation.

\begin{proposition}\label{prop:strict}
For $h(x)=x$,
\[
 \|h\|_{\bmo(\T)}=1,
 \qquad
 \|h\|_{\BMO(\T)}=\infty.
\]
In particular, on centered $L_2(\g)$ elements one has a strict inclusion
$\BMO(\T)\subsetneq\bmo(\T)$.
\end{proposition}

\begin{proof}
Mehler's formula gives
\[
 T_th=e^{-t}h,
 \qquad
 T_th^2=e^{-2t}h^2+1-e^{-2t}.
\]
The first assertion follows immediately from \eqref{eq:bmo}.  On the other
hand,
\[
 T_t|h-T_th|^2
 =(1-e^{-t})^2\big(e^{-2t}h^2+1-e^{-2t}\big),
\]
which is unbounded for every $t>0$.  The source proves under $\Ga_2\geq0$
that $\bmo\lesssim\BMO$ \cite[Proposition 2.4]{JungeMei}, so the inclusion
holds; the coordinate makes it strict.
\end{proof}

\section{Separation of the norms in part (a)}

Let $P_s=e^{-s\sqrt A}$ be the Poisson semigroup subordinated to $T$.  The
source defines
\begin{align*}
 \|f\|_{\BMO(\Ga)}^2
 &=\sup_{t>0}\left\|P_t\int_0^t\Ga(P_sf)\,s\,ds\right\|_\infty,\\
 \|f\|_{\BMO(\partial)}^2
 &=\sup_{t>0}\left\|P_t\int_0^t|\partial_sP_sf|^2s\,ds\right\|_\infty.
\end{align*}

\begin{theorem}\label{thm:poisson}
For $h(x)=x$,
\[
 \|h\|_{\BMO(\Ga)}=\frac12,
 \qquad
 \|h\|_{\BMO(\Pp)}=\|h\|_{\BMO(\partial)}=\infty.
\]
Furthermore,
\[
 \sup_R\|f_R\|_{\BMO(\Ga)}\leq\frac12,
 \qquad
 \|f_R\|_{\BMO(\Pp)}\longrightarrow\infty.
\]
Thus symmetry, ergodicity, analyticity, hypercontractivity and
$\Ga_2\geq\Ga$ do not suffice for the triple equivalence in open problem
(ii)(a), and the failure is visible through bounded approximants.
\end{theorem}

\begin{proof}
The coordinate lies in the first Wiener chaos: $Ah=h$.  Also
$A(h^2-1)=2(h^2-1)$.  Hence
\begin{equation}\label{eq:poisson-eigs}
 P_sh=e^{-s}h,
 \qquad
 P_th^2=1+e^{-\sqrt2t}(h^2-1).
\end{equation}
Since $\Ga(P_sh)=e^{-2s}$,
\[
 \|h\|_{\BMO(\Ga)}^2
 =\sup_{t>0}\int_0^t e^{-2s}s\,ds=\frac14.
\]
But
\[
 P_t|h-P_th|^2=(1-e^{-t})^2P_th^2
\]
is unbounded by \eqref{eq:poisson-eigs}; this is the $\BMO(\Pp)$ assertion.
Similarly,
\[
 P_t\int_0^t|\partial_sP_sh|^2s\,ds
 =\left(\int_0^t e^{-2s}s\,ds\right)P_th^2
\]
is unbounded, proving the $\BMO(\partial)$ assertion.

For the bounded truncations, subordination and the Mehler derivative identity
give
\[
 |(P_sf_R)'|
 \leq\int_0^\infty\eta_s(u)|(T_uf_R)'|\,du
 \leq\int_0^\infty\eta_s(u)e^{-u}\,du=e^{-s},
\]
where $\eta_s$ is the positive subordination density.  Therefore
$\Ga(P_sf_R)\leq e^{-2s}$ and
$\|f_R\|_{\BMO(\Ga)}\leq1/2$.

Finally $f_R\to h$ in $L_2(\g)$ and
$(P_t-P_{2t})h=(e^{-t}-e^{-2t})h$.  Repeating the $L_2$-limit argument in
\eqref{eq:L2limit} shows
$\|(P_t-P_{2t})f_R\|_\infty\to\infty$ for every fixed $t>0$.  Kadison's
inequality for $P_t$ yields
\[
 \|(P_t-P_{2t})f_R\|_\infty\leq\|f_R\|_{\BMO(\Pp)},
\]
and completes the proof.
\end{proof}

\section{Verification, scope, and novelty bounds}

The proof uses only Mehler's formula, Gaussian $L_2$ convergence, positive
subordination, and Kadison's inequality.  No computation is needed.  The
included script \texttt{code/verify\_ou\_truncations.py} evaluates the first
two moments of a clipped Gaussian exactly in terms of the normal cdf and
checks on a finite grid that the transition variance is at most one while the
fixed-time drift grows linearly with $R$.  This is a diagnostic, not evidence
for the proof.

The result fully refutes the universal estimates (c) and (c').  It also shows
that the strong hypotheses listed in Theorem \ref{thm:poisson} are
insufficient for (a).  The broad request for a complete characterization of
all semigroups remains open.  Part (b) is not challenged: the source already
proves its Poisson comparison under $\Ga_2\geq0$.

A bounded novelty search on 9 August 2026 covered the run's registry,
solution, attempt, and proof-gap indexes; exact and close-variant arXiv
searches combining the displayed norm comparison, the $T_t-T_{2t}$ term, and
Ornstein--Uhlenbeck; the original paper; Mei's follow-up
\cite{MeiDuality}; the Poisson $H^\infty$ paper \cite{FergusonMeiSimanek};
and the later algebraic Calderón--Zygmund treatment of Gaussian semigroup BMO
\cite{JungeMeiParcetXia}.  The follow-up explicitly says the two norms may
differ, proves equality under extra assumptions, and notes that OU fails one
of those assumptions, but the searched sources did not contain the explicit
coordinate/truncation separation above or the part-(a) separation.  Novelty
is therefore provisional; this was not an exhaustive MathSciNet or zbMATH
priority search.

\textbf{Human-review recommendation.} High priority.  Verify the source's
domain conventions, the Lipschitz conditional-variance estimate, and the
subordination gradient bound.  The bounded sequence $f_R$ makes the failure
of (c),(c') independent of whether the source's completed BMO spaces are
defined from $L_\infty$ or $L_2$.

\sloppy
\begin{thebibliography}{9}
\raggedright
\bibitem{JungeMei}
M. Junge and T. Mei,
\emph{BMO spaces associated with semigroups of operators},
Math. Ann. 352 (2012), 691--743; arXiv:1111.6194.
\url{https://arxiv.org/abs/1111.6194}.

\bibitem{MeiDuality}
T. Mei,
\emph{An $H^1$--BMO duality theory for semigroups of operators},
arXiv:1204.5082 (2012).
\url{https://arxiv.org/abs/1204.5082}.

\bibitem{FergusonMeiSimanek}
T. Ferguson, T. Mei, and B. Simanek,
\emph{$H^\infty$-calculus for semigroup generators on BMO},
Adv. Math. 347 (2019), 408--441; arXiv:1701.06623.
\url{https://arxiv.org/abs/1701.06623}.

\bibitem{JungeMeiParcetXia}
M. Junge, T. Mei, J. Parcet, and R. Xia,
\emph{Algebraic Calderón--Zygmund theory},
Adv. Math. 376 (2021), 107443; arXiv:1907.07375.
\url{https://arxiv.org/abs/1907.07375}.
\end{thebibliography}

\end{document}
